\chapter{Planos base u originales}

\section{Objeto}
Este capítulo documenta los planos originales recibidos como fuente del
proyecto y explica como sirvieron de base para el desarrollo del anteproyecto
electrico. Los archivos DWG originales se conservan intactos en
\texttt{fuentes/local/cad/} y no se modifican. La geometria de trabajo se
extrajo de ellos y se estructuro en
\texttt{arquitectura/datos/layout-grifo.json}.

\section{Registro de planos originales}

\begin{center}
\small
\begin{tabularx}{\linewidth}{@{}l p{3.4cm} p{3.4cm} X@{}}
\toprule
Campo & CAD-001 (ubicacion) & CAD-002 (distribucion) & Nota \\
\midrule
Nombre del plano & PLANO DE UBICACION & PLANO DE DISTRIBUCION &
  nombre del archivo DWG recibido \\
Codigo & CAD-001 & CAD-002 & codigo interno del proyecto \\
Escala & no verificado en fuente & no verificado en fuente &
  verificar el cajetin del DWG original \\
Fecha & no disponible en fuente & no disponible en fuente &
  consultar sello del plano original \\
Version & unica recibida & unica recibida & \\
Descripcion & Ubicacion del predio Rustico Reumita, Parcela B-8 y B-9, comunidad
  San Francisco de Buenavista, carretera Juliaca--Puno & Planta arquitectonica
  de distribucion: ambientes, tanques, surtidores, circulacion y equipos
  electricos observados &
  texto transcrito por el estudiante del rotulo \\
Observaciones & Distrito y provincia San Roman, departamento Puno; altitud
  aproximada 3830~m s.n.m. & Se observan TG, TG2, PAT, pararrayo (R=20~m,
  h=12~m), pulsador de emergencia, tanques TK-1..TK-3 y puntos PM (ambiguos,
  por confirmar) &
  CAD-002 es la fuente de la geometria de trabajo \\
\bottomrule
\end{tabularx}
\end{center}

\section{Huella de verificacion de las fuentes}
La integridad de los archivos originales se registra con su resumen SHA-256 en
\texttt{proyecto.yaml}:
\begin{itemize}
\item CAD-001 (PLANO DE UBICACION.dwg): \texttt{35d41698d6\-4d82adec580\-cd0cf17424ea\-6fa55ec3dd9702\-ace8ab112566d40ae}.
\item CAD-002 (PLANO DE DISTRIBUCION.dwg): \texttt{2ddd45113a7d94\-eb27b595272c3f\-e6d9ea2142999354\-b5cf4c350247f3db506b}.
\end{itemize}

\section{Uso de cada plano original}
\subsection{CAD-001 (ubicacion)}
Sirvio para confirmar la ubicacion geografica del predio, el distrito y la
provincia (San Roman, Puno) y la condicion rural del lote. Estos datos definen
la altitud de referencia (3830~m s.n.m.), que condiciona el desrating de los
conductores y la correccion por altitud del grupo electrogeno, asi como la
empresa distribuidora asumida (Electro Puno).

\subsection{CAD-002 (distribucion)}
Es la fuente principal de la geometria. De el se extrajeron:
\begin{itemize}
\item Los ambientes (administracion, oficina 01, sala de maquinas, SS.HH),
  cuyas dimensiones alimentan la memoria de calculo de iluminacion.
\item La zona de tanques (TK-1 Gasohol Premium, TK-2 Gasohol Regular, TK-3
  Diesel B5-S50) y las islas de despacho, que definen la ubicacion de las
  bombas sumergibles (STP), las cabezas electronicas y las areas peligrosas de
  la lamina IE-06.
\item Los equipos electricos observados (TG, TG2, interruptor general, PAT,
  pararrayo, pulsador de emergencia), que se replantearon como tableros TDE,
  TDF y TD-A1 en el diseno del anteproyecto.
\end{itemize}

\section{Reconstruccion vectorial de la geometria}
Debido a que no se dispone de un convertidor de DWG en el entorno de trabajo,
la vista esquematica de la planta siguiente se reconstruyo a partir de la
geometria extraida y validada (\texttt{layout-grifo.json}), y es la base sobre
la que se dibujaron las laminas IE-01 a IE-06.

\IlustracionLayout

\begin{center}
\small
\begin{tabularx}{\linewidth}{@{}l X@{}}
\toprule
Elemento & Uso en el diseno electrico \\
\midrule
Ambientes administrativos & definen los circuitos de alumbrado y
  tomacorrientes de TD-A1 y los calculos de iluminacion \\
Zona de tanques e islas & definen las STP, cabezas electronicas y los limites
  academicos de zonas 1 y 2 (IE-06) \\
TG / TG2 / interruptor general & replanteados como tablero general y tableros
  de emergencia, fuerza y administracion \\
PAT, pararrayo y pulsador & definen la puesta a tierra, la proteccion contra
  el rayo y el paro de emergencia (IE-03, IE-04) \\
\bottomrule
\end{tabularx}
\end{center}
